% !TEX root = ../main.tex
% =====================================================================
% 结论（附件 3 目录中为“5  结论（黑体四号）”）
%
% 以下文字与附件 3 第 5 章完全一致，逐字照录（含附件 3 自带的格式说明
% 段落）。章标题照录为“5  结论”。
% =====================================================================

\ujschapter{5\quad 结论}

\section{***（宋体四号，加粗，1.5倍行距，段前段后0.5行）}

结论应集中反映作者的研究成果，要求语言精炼、表述准确且完整，着重阐述自己的创造性成果及其在本研究领域中的意义、作用，同时应包括所得结果与已有结果的比较和本课题尚存在的问题，以及进一步开展研究的见解与建议。

\subsection{***（宋体小四，加粗，1.5倍行距，段前段后0.5行）}

……………………（正文小四，中文宋体，英文数字Times New Roman，1.5行距，首行缩进2字符）

\section{***（宋体四号，加粗，1.5倍行距，段前段后0.5行）}

\subsection{***（宋体小四，加粗，1.5倍行距，段前段后0.5行）}

……………………（正文小四，中文宋体，英文数字Times New Roman，1.5行距，首行缩进2字符）
